\documentclass[10pt]{article}

\usepackage[T1]{fontenc}
\usepackage[margin=0.55in]{geometry}
\usepackage{amsmath}
\usepackage{array}
\usepackage{booktabs}
\usepackage{enumitem}
\usepackage{longtable}
\usepackage{pdflscape}
\usepackage{ragged2e}
\usepackage{xurl}

\newcolumntype{P}[1]{>{\RaggedRight\arraybackslash}p{#1}}
\newcommand{\celllist}[1]{%
  \begin{minipage}[t]{\linewidth}
    \vspace{0pt}
    \begin{itemize}[leftmargin=1.1em, itemsep=2pt, topsep=0pt, parsep=0pt]
      #1
    \end{itemize}
  \end{minipage}%
}

\begin{document}

\begin{landscape}
\small
\setlength{\tabcolsep}{4pt}
\renewcommand{\arraystretch}{1.15}

\begin{longtable}{@{}P{0.10\linewidth} P{0.20\linewidth}
                         P{0.36\linewidth} P{0.27\linewidth}@{}}
  \caption{Formalization status of the Kleinberg--Tardos case studies. ``Proved'' means
  theorem-checked in Lean from the repository's executable or interpreted model. The final
  column records missing end-to-end connections, stated restrictions, and natural extensions;
  it does not weaken the theorems already listed.}
  \label{tab:case-study-status} \\

  \toprule
  \textbf{Case study} &
  \textbf{Textbook results targeted} &
  \textbf{What is proved in Lean} &
  \textbf{What remains open or restricted} \\
  \midrule
  \endfirsthead

  \multicolumn{4}{c}{\tablename\ \thetable{} --- continued} \\
  \toprule
  \textbf{Case study} &
  \textbf{Textbook results targeted} &
  \textbf{What is proved in Lean} &
  \textbf{What remains open or restricted} \\
  \midrule
  \endhead

  \midrule
  \multicolumn{4}{r}{Continued on the next page} \\
  \endfoot

  \bottomrule
  \endlastfoot

  \textbf{BFS} &
  \celllist{
    \item Layer-based reachability and shortest-path-tree correctness.
    \item Theorem~(3.11): adjacency-list time $O(m+n)$.
    \item Theorem~(3.15): BFS layers either give a proper bipartite coloring or expose an odd
      cycle.
  } &
  \celllist{
    \item \textbf{Execution-linked:} the completed Boolean-table run visits exactly the vertices
      reachable from the source (\path{run_visitsExactlyReachable}).
    \item For every finite connected graph, a certified spanning tree contained in the graph is
      constructed and proved to preserve all source distances.
    \item Theorem~(3.15) is proved from that shortest-path-tree certificate
      (\path{bfs_bipartite_or_contains_odd_cycle}).
    \item Exact measured-run bounds are proved for arbitrary bounded primitive costs. They yield
      $O(m+n)$ for adjacency lists and $O(n^2)$ for adjacency matrices
      (\path{adjacencyListTime_isBigO}, \path{adjacencyMatrixTime_isBigO}).
  } &
  \celllist{
    \item Prove that the concrete edge list and level table emitted by the executable run are
      themselves the certified shortest-path tree and distance layers.
    \item Connect the executable output directly to the coloring/odd-cycle witnesses in
      Theorem~(3.15); the current certificate is a separate noncomputable construction.
  } \\
  \midrule

  \textbf{DFS} &
  \celllist{
    \item Reachability of source-based depth-first search.
    \item Theorem~(3.7): every non-tree edge of an undirected DFS tree joins an
      ancestor--descendant pair.
    \item Theorem~(3.13): adjacency-list time $O(m+n)$.
  } &
  \celllist{
    \item \textbf{Execution-linked:} the iterative Boolean-table run visits exactly the vertices
      reachable from the source (\path{run_visitsExactlyReachable}).
    \item Theorem~(3.7) is proved from an abstract recursive-DFS tree certificate with discovery
      and return times (\path{RecursiveTreeCertificate.nonTreeEdge_ancestor}).
    \item Exact measured-run bounds are proved for arbitrary bounded primitive costs, giving
      $O(m+n)$ for adjacency lists and $O(n^2)$ for adjacency matrices
      (\path{adjacencyListTime_isBigO}, \path{adjacencyMatrixTime_isBigO}).
  } &
  \celllist{
    \item Construct the recursive-tree certificate, including discovery/return times and tree
      edges, from the actual iterative DFS execution.
    \item This connection would make Theorem~(3.7) end-to-end for the executable runner rather
      than conditional on a separately supplied certificate.
  } \\
  \midrule

  \textbf{Mergesort} &
  \celllist{
    \item Sorted-permutation correctness.
    \item Theorem~5.1: the recurrence $T(n) \leq 2T(n/2)+cn$.
    \item Theorem~5.2: substitution gives $T(n)=O(n\log n)$.
  } &
  \celllist{
    \item The pure program and every measured interpretation return a sorted permutation
      (\path{mergeSort_correct}, \path{interpret_mergeSort_correct}).
    \item A cost-independent operation-profile proof bounds comparison, split, merge, and
      base-case work; any cost model satisfying the primitive bounds inherits the result.
    \item Theorem~5.1 is formalized as a recurrence upper bound, and Theorem~5.2 is proved by
      substitution: $T(2^k) \leq c\,2^k k$
      (\path{theorem_5_2_substitution}).
    \item Exact measured cost is proved $O(n\log n)$ for every power-of-two input family
      (\path{mergeSort_exactCost_isBigO_on_powersOfTwo}).
  } &
  \celllist{
    \item Remove the $n=2^k$ restriction by solving the executable floor/ceiling recurrence for
      arbitrary input lengths.
    \item Refine aggregate structural charges into explicit costs for operations such as
      \path{length}, \path{take}, \path{drop}, allocation, or peak live space if those metrics are
      desired.
  } \\
  \midrule

  \textbf{Dijkstra} &
  \celllist{
    \item Claim~(4.14): Dijkstra settles vertices with correct shortest-path distances.
    \item Claim~(4.15): with adjacency lists and a heap, the running time is
      $O(m\log n)$.
  } &
  \celllist{
    \item \textbf{Execution-linked, under reachability:} for nonnegative weights, the actual
      run's predecessor array reconstructs exact edge-occurrence paths of the stored weight, and
      every settled vertex receives a shortest path (\path{claim_4_14}); the result is lifted to
      both measured runners.
    \item The actual trace has exactly $n$ extract-min operations, exactly $m$ row-scan and $m$
      relaxation-inspection units, and at most $m$ key changes
      (\path{claim_4_15_adjacencyList}).
    \item The logarithmic-heap run satisfies an explicit finite bound, then
      $O((m+n)\log(n+1))$; for $n\geq2$ and reachable inputs, this yields the textbook
      $O(m\log n)$ corollary (\path{claim_4_15_heap_corollary}).
    \item The same development covers \path{Nat}, \path{NNRat}, and \path{NNReal} weights through
      one generic nonnegative-weight interface.
  } &
  \celllist{
    \item Remove the assumption that every enumerated vertex is reachable, and specify/prove the
      final infinite distances and predecessor entries for disconnected graphs.
    \item Turn the existing source-level space components into a formal space or peak-live-space
      theorem; no such theorem is currently claimed.
  } \\
  \midrule

  \textbf{Randomized quicksort} &
  \celllist{
    \item Sorted-permutation correctness for randomized pivot choices.
    \item Section~13.5's unnumbered expected-time claim $O(n\log n)$.
    \item Worst-case running time $\Theta(n^2)$.
  } &
  \celllist{
    \item Every supported result of the actual uniform-pivot runner is a sorted permutation
      (\path{quicksort_result_correct}).
    \item Exact comparison and unit-textbook expected-cost recurrences are derived from the joint
      probability--cost semantics.
    \item On distinct inputs, the expected comparison cost and the unit textbook cost are proved
      $O(n\log n)$ (\path{expectedComparisonsBySize_isBigO},
      \path{expectedTextbookRuntimeBySize_isBigO}).
    \item The maximum supported comparison cost is exactly $\binom{n}{2}$, and both comparison
      cost and textbook cost have proved $\Theta(n^2)$ worst-case bounds
      (\path{worstComparisonsBySize_isTheta}, \path{worstTextbookRuntimeBySize_isTheta}).
  } &
  \celllist{
    \item Generalize the expected- and worst-cost analyses to duplicate keys.
    \item Prove an exact closed form for the expected number of comparisons, rather than only a
      harmonic upper envelope.
    \item Formalize the modified central-splitter rejection algorithm and the section's Select
      algorithm.
    \item Add an in-place array implementation or an allocation/peak-space analysis; neither is
      part of the current model.
  } \\

\end{longtable}
\end{landscape}

\end{document}
